% !TeX spellcheck = en

%%%%%%%%%%%%%%%%%%%%%%%%%%%%%%%%%%%%%%%%%%%%%%%%%%%%%%%%%%%%%%%%%%%%%%%%%%%%%%%%%%
% SECTION
%%%%%%%%%%%%%%%%%%%%%%%%%%%%%%%%%%%%%%%%%%%%%%%%%%%%%%%%%%%%%%%%%%%%%%%%%%%%%%%%%%	
\section{Beispiel Abschnitt}
\label{sec:some}
%\setcounter{subsection}{1}

%%%%%%%%%%%%%%%%%%%%%%%%%%%%%%%%%%%%%%%%%%%%%%%%%%%%%%%%%%%%%%%%%%%%%%%%%%%%%%%%%%
% FRAME
%%%%%%%%%%%%%%%%%%%%%%%%%%%%%%%%%%%%%%%%%%%%%%%%%%%%%%%%%%%%%%%%%%%%%%%%%%%%%%%%%%
\begin{frame}{Eine Beispiel Folie}
	
Inhalt
	
\end{frame}


\section{Grundlagen}

\subsection{README}

\begin{frame}[fragile]
	% [fragile] braucht man, wenn man \verb oder \verbatim auf der Folie verwendet!
	\frametitle{README}
	\begin{itemize}
		\item Das hier ist die Folienvorlage zum Kurs \textbf{Betriebssysteme und Rechnernetze}.
		\item Diese Vorlage soll den Einstieg erleichtern.
		\item Wird unter Linux im Verzeichnis mit der Quelldatei (\texttt{.tex}) das Kommando \texttt{make} eingegeben, wird eine PDF- und eine PS-Datei erzeugt. Beide haben den gleichen Inhalt.
		\item Ein aktuelles \LaTeX\ sollte installiert sein.
		\item Als Editor kann ein Texteditor (bspw. Kate) oder eine IDE wie bspw. TeXStudio \cite{TeXStudio} verwendet werden.
		\item Wer unter Windows die Folien machen möchte, dem empfehle ich MiK\TeX\footnote{\texttt{http://www.miktex.org}} \cite{MikTeX} und \TeX nicCenter\footnote{\texttt{http://www.texniccenter.org}} \cite{TexNic}. Hierzu habe ich aber keine Erfahrungswerte.
	\end{itemize}
\end{frame}

\begin{frame}[fragile]
	\frametitle{\LaTeX Beamer}
	\begin{itemize}
		\item Diese Vorlage nutzt die \LaTeX-Klasse \texttt{beamer}.
		\item Eine gute Dokumentation über diese Klasse befindet sich hier:\\
		\url{http://www2.informatik.hu-berlin.de/~mischulz/beamer.html}
		\item Diese Quelle ist auch hilfreich:\\
		\url{http://www.physik.uni-freiburg.de/~tooleh/latex_beamerkurs.pdf}
		\item Google findet sehr viele hilfreiche Links zum Thema \textbf{LaTeX Beamer}.
	\end{itemize}
\end{frame}

\section{Textformatierung}

\subsection{Schriften und Sonderzeichen}

\begin{frame}[fragile]
	\frametitle{Schriften und Sonderzeichen}
	\begin{itemize}
		\item Es gibt verschiedene Schriftsätze: \textbf{Bild Face}, \textrm{Roman}, \textit{Italic}, \texttt{Typewriter}, \textsf{Sans Serif}, \textsl{Slanted}, \textsc{Small Caps}.
		\item \textcolor{blue}{Farben} sollte man \textcolor{red}{nicht} zu viel einsetzen.
		\item Ein paar Sonderzeichen: \textbackslash, \$, \&, \euro, \%, \#, \textunderscore, \textasciitilde, \textasciicircum, \textbar, \{, \}
		\item Weitere Sonderzeichen: \copyright, \textregistered, \texttrademark, \S, \P, \pounds, \dag, \ddag, \textbullet
		\item Fortsetzungspunkte macht das Kommando \verb!\dots!. Ergebnis: \dots
	\end{itemize}
\end{frame}

\subsection{Schriftgrößen}

\begin{frame}[fragile]
	\frametitle{Schriftgrößen}
	
	\Huge
	\verb!\Huge!
	\normalsize
	
	\huge
	\verb!\huge!
	\normalsize
	
	\LARGE
	\verb!\LARGE!
	\normalsize
	
	\Large
	\verb!\Large!
	\normalsize
	
	\large
	\verb!\large!
	\normalsize
	
	\normalsize
	\verb!\normalsize!
	\normalsize
	
	\small
	\verb!\small!
	\normalsize
	
	\footnotesize
	\verb!\footnotesize!
	\normalsize
	
	\scriptsize
	\verb!\scriptsize!
	\normalsize
	
	\tiny
	\verb!\tiny!
	\normalsize
\end{frame}

\subsection{Blöcke}

\begin{frame}[t]
	% [t] legt fest, dass der Text oben auf der Folie platziert wird.
	\frametitle{Blöcke}
	\begin{itemize}
		\item Es gibt verschiedene Arten von Blöcken:
	\end{itemize}
	\begin{block}{Blocktitel}
		Blocktext
	\end{block}
	
	\begin{exampleblock}{Blocktitel}
		Blocktext
	\end{exampleblock}
	
	\begin{alertblock}{Blocktitel}
		Blocktext
	\end{alertblock}
\end{frame}

\subsection{Bilder}

\begin{frame}
	% [t] legt fest, dass der Text oben auf der Folie platziert wird.
	\frametitle{Bilder}
	\begin{itemize}
		\item Hier ist ein Bild:
	\end{itemize}
	\begin{center}
		\includegraphics[height=4cm]{Pinguin.pdf}
		% Die Dateiendung muss man nicht angeben.
		% Man hätte auch alternativ die Breite mit "width" angeben können.
	\end{center}
	\begin{itemize}
		\item Bilder sollten im Format Encapsulated PostScript (\texttt{.eps}) sein. Dieses Dateiformat kann man mit Gimp und vielen anderen Programmen erzeugen.
	\end{itemize}
\end{frame}

\subsection{Tabellen}

\begin{frame}[fragile]
	\frametitle{Tabellen}
	\begin{itemize}
		\item Es gibt mehrere Umgebungen, um Tabellen zu machen. \texttt{tabular} ist nur eine von vielen.
	\end{itemize}
	
	\begin{center}
		\begin{tabular}{|c|l|c|r|}
			\hline
			\textbf{Zeile} & \textbf{Linksbündig} & \textbf{Zentriert} & \textbf{Rechtsbündig} \\
			\hline\hline
			1 & Zeile 1 & Zeile 1 & Zeile 1 \\
			2 & Zeile 2 & Zeile 2 & Zeile 2 \\
			3 & Zeile 3 & Zeile 3 & Zeile 3 \\
			\hline
		\end{tabular}
	\end{center}
	
	\begin{itemize}
		\item Das geht natürlich auch ohne die Rahmen:
	\end{itemize}
	
	\begin{center}
		\begin{tabular}{clcr}
			\textbf{Zeile} & \textbf{Linksbündig} & \textbf{Zentriert} & \textbf{Rechtsbündig} \\
			1 & Zeile 1 & Zeile 1 & Zeile 1 \\
			2 & Zeile 2 & Zeile 2 & Zeile 2 \\
			3 & Zeile 3 & Zeile 3 & Zeile 3 \\
		\end{tabular}
	\end{center}
\end{frame}

\subsection{Mehrspaltige Folien}

\begin{frame}
	\frametitle{Mehrspaltige Folien}
	\begin{columns}
		\column{.45\textwidth}
		Mehrspaltige Folien können einfach mit \texttt{columns} realisiert werden.
		\column{.45\textwidth}
		\begin{enumerate}
			\item Ein Eintrag
			\item Noch ein Eintrag
		\end{enumerate}
	\end{columns}
\end{frame}